% Tiểu luận chuyên ngành: RethinkMCTS trên các LLM mở cỡ nhỏ (7–9B tham số).
% Biên dịch: make   (xelatex → biber → xelatex ×2). Quy cách: tlcn.sty; thông tin bìa: thongtin.tex.
\documentclass[fontsize=13pt,paper=a4,oneside,open=any,numbers=noenddot,listof=totoc,
  bibliography=totoc,headings=normal]{scrreprt}
\usepackage{tlcn}
\addbibresource{tailieu.bib}
% Thông tin chung dùng cho bìa, bìa phụ, phiếu nhận xét.
% Trường mới đổi tên từ ĐH Sư phạm Kỹ thuật TP.HCM thành ĐH Công nghệ Kỹ thuật TP.HCM.
\newcommand{\truong}{TRƯỜNG ĐẠI HỌC CÔNG NGHỆ KỸ THUẬT TP.HCM}
\newcommand{\truongngan}{ĐH CÔNG NGHỆ KỸ THUẬT TP.HCM}
\newcommand{\khoa}{KHOA CÔNG NGHỆ THÔNG TIN}
\newcommand{\bomon}{BỘ MÔN TRÍ TUỆ NHÂN TẠO}
\newcommand{\nganh}{Công nghệ Thông tin}

\newcommand{\sinhvienmot}{Phan Quốc Viễn}
\newcommand{\mssvmot}{23110362}
\newcommand{\sinhvienhai}{Nghiêm Quang Huy}
\newcommand{\mssvhai}{23110222}

% Tên đề tài (giữ đúng tên trong Proposal.md).
\newcommand{\tendetai}{Cải thiện khả năng suy luận của LLM trong sinh mã nguồn bằng Monte Carlo Tree Search với phản hồi từ kiểm thử tự động}
\newcommand{\TENDETAI}{CẢI THIỆN KHẢ NĂNG SUY LUẬN CỦA LLM TRONG SINH MÃ NGUỒN BẰNG MONTE CARLO TREE SEARCH VỚI PHẢN HỒI TỪ KIỂM THỬ TỰ ĐỘNG}

\newcommand{\loaikhoaluan}{TIỂU LUẬN CHUYÊN NGÀNH}
\newcommand{\gvhd}{TS. Phan Thị Huyền Trang}
\newcommand{\khoahoc}{KHÓA 2023 -- 2027}
\newcommand{\namthuchien}{2026}


% Ô giữ chỗ cho phần chưa viết; xóa khi viết xong (tìm "lipsumtodo").
\newcommand{\lipsumtodo}[1]{\par\noindent\fcolorbox{orange}{orange!8}{\parbox{\dimexpr\linewidth-2\fboxsep-2\fboxrule}{\setstretch{1.1}\small\textbf{Cần viết:} #1}}\par}

% Bìa chính (có khung) và bìa phụ (không khung), theo mẫu trong mau/Mau-5-TrinhBay-KLTN.docx.
% \trangbia{true} = bìa chính có khung đôi; \trangbia{false} = bìa phụ.
\newcommand{\trangbia}[1]{%
  \begin{titlepage}
    \ifstrequal{#1}{true}{%
      \begin{tikzpicture}[remember picture,overlay]
        \draw[line width=1.6pt] ([xshift=2.4cm,yshift=-1.4cm]current page.north west) rectangle ([xshift=-1.4cm,yshift=1.4cm]current page.south east);
        \draw[line width=0.5pt] ([xshift=2.55cm,yshift=-1.55cm]current page.north west) rectangle ([xshift=-1.55cm,yshift=1.55cm]current page.south east);
      \end{tikzpicture}\par}{}%
    \setstretch{1.2}\setlength{\parindent}{0pt}
    \centering
    {\fontsize{14pt}{18pt}\selectfont\bfseries \truong\par \khoa\par \bomon\par}
    \vspace{1.2cm}
    % Logo khoa CNTT (FIT HCM-UTE) theo mẫu, do người dùng cung cấp. Logo trường (hinh/logo-truong.png) để dự phòng.
    \IfFileExists{hinh/logo-cntt.png}{\includegraphics[height=3.2cm]{logo-cntt.png}}{\fbox{\parbox[c][3cm][c]{3cm}{\centering\small logo khoa CNTT\\(hinh/logo-cntt.png)}}}\par
    \vspace{1.2cm}
    {\fontsize{14pt}{18pt}\selectfont\bfseries \MakeUppercase{\sinhvienmot} -- \mssvmot\par \MakeUppercase{\sinhvienhai} -- \mssvhai\par}
    \vspace{1.5cm}
    {\fontsize{14pt}{18pt}\selectfont Đề tài:\par}
    \vspace{0.3cm}
    {\fontsize{16pt}{22pt}\selectfont\bfseries \TENDETAI\par}
    \vspace{1.5cm}
    {\fontsize{14pt}{18pt}\selectfont\bfseries \loaikhoaluan\par}
    \vspace{1.2cm}
    {\fontsize{13pt}{17pt}\selectfont GIẢNG VIÊN HƯỚNG DẪN\par \bfseries \MakeUppercase{\gvhd}\par}
    \vfill
    {\fontsize{13pt}{17pt}\selectfont\bfseries \khoahoc\par}
    \vspace{0.6cm}
  \end{titlepage}}

% Phiếu nhận xét của GVHD / GVPB, theo mẫu. \phieunhanxet{HƯỚNG DẪN}{hướng dẫn}{tên GV, để trống thì in dấu chấm}
\newcommand{\phieunhanxet}[3]{%
  \clearpage\thispagestyle{empty}
  {\setstretch{1.2}\fontsize{12pt}{14pt}\selectfont\noindent
  \begin{tabular}{@{}p{0.4\linewidth}@{}p{0.6\linewidth}@{}}
    \centering\bfseries \truongngan\\ \centering\bfseries KHOA CNTT\\ \centering *******
    & \centering\bfseries CỘNG HOÀ XÃ HỘI CHỦ NGHĨA VIỆT NAM\\ \centering\bfseries Độc lập -- Tự do -- Hạnh phúc\\ \centering *******
  \end{tabular}\par}
  \vspace{1cm}
  {\centering\fontsize{14pt}{18pt}\selectfont\bfseries PHIẾU NHẬN XÉT CỦA GIẢNG VIÊN #1\par}
  \vspace{0.6cm}
  \noindent Họ và tên Sinh viên 1: \sinhvienmot \hfill MSSV 1: \mssvmot\par
  \noindent Họ và tên Sinh viên 2: \sinhvienhai \hfill MSSV 2: \mssvhai\par
  \noindent Ngành: \nganh\par
  \noindent Tên đề tài: \tendetai\par
  \noindent Họ và tên Giảng viên #2: \ifstrempty{#3}{\dotfill}{#3}\par
  \vspace{0.3cm}
  {\centering\bfseries NHẬN XÉT\par}
  \begin{enumerate}[label=\arabic*.,leftmargin=*]
    \item Về nội dung đề tài \& khối lượng thực hiện:\\ \dotfill\\ \dotfill\\ \dotfill
    \item Ưu điểm:\\ \dotfill\\ \dotfill
    \item Khuyết điểm:\\ \dotfill\\ \dotfill
    \item Đề nghị cho bảo vệ hay không?\\ \dotfill
    \item Đánh giá loại: \dotfill
    \item Điểm: \dotfill
  \end{enumerate}
  \vspace{0.5cm}
  \hfill\begin{minipage}{0.6\linewidth}\centering\setstretch{1.2}
    \emph{Tp. Hồ Chí Minh, ngày \quad tháng \quad năm \namthuchien}\\
    \textbf{Giảng viên #2}\\ (Ký \& ghi rõ họ tên)
  \end{minipage}
  \clearpage}


\begin{document}

% ------------------------------------------------------------------ Phần đầu (không đánh số / số La Mã)
\trangbia{true}   % bìa chính
\trangbia{false}  % bìa phụ
\phieunhanxet{HƯỚNG DẪN}{hướng dẫn}{\gvhd}
\phieunhanxet{PHẢN BIỆN}{phản biện}{}

\pagenumbering{roman}
\addchap{LỜI CẢM ƠN}
% TODO: viết lời cảm ơn (GVHD, khoa, gia đình...).
\lipsumtodo{Lời cảm ơn.}

% Đề cương chi tiết có chữ ký của GVHD (chấp nhận bản photo/scan).
% Đặt bản scan PDF vào decuong/decuong-ky.pdf; nếu chưa có, in một trang giữ chỗ.
\IfFileExists{decuong/decuong-ky.pdf}{%
  \includepdf[pages=-,pagecommand={\thispagestyle{empty}},addtotoc={1,chapter,0,ĐỀ CƯƠNG CHI TIẾT,decuong}]{decuong/decuong-ky.pdf}}{%
  \addchap{ĐỀ CƯƠNG CHI TIẾT}
  \lipsumtodo{Chèn bản scan đề cương chi tiết có chữ ký của GVHD vào \texttt{decuong/decuong-ky.pdf}.}}

\tableofcontents
\listoffigures
\listoftables
\addchap{DANH MỤC KÝ HIỆU VÀ CHỮ VIẾT TẮT}
% Xếp theo thứ tự bảng chữ cái. Chỉ ghi các từ viết tắt thực sự dùng nhiều lần trong tiểu luận.
{\setstretch{1.2}
\begin{longtable}{@{}p{0.22\linewidth}p{0.74\linewidth}@{}}
\toprule
\textbf{Viết tắt} & \textbf{Ý nghĩa} \\
\midrule
\endhead
\bottomrule
\endfoot
APPS & Automated Programming Progress Standard (bộ dữ liệu sinh mã) \\
BoN & Best-of-$N$: sinh $N$ chương trình độc lập, chọn theo kết quả public test \\
CFG & Control Flow Graph (đồ thị luồng điều khiển) \\
CI & Confidence Interval (khoảng tin cậy) \\
LLM & Large Language Model (mô hình ngôn ngữ lớn) \\
MCTS & Monte Carlo Tree Search \\
P-UCB & Predictor Upper Confidence Bound (công thức chọn nút trong MCTS) \\
\end{longtable}}

\addchap{TÓM TẮT TIỂU LUẬN}
% Tóm tắt vấn đề nghiên cứu, hướng tiếp cận, cách giải quyết và kết quả chính (mọi con số lấy từ
% docs/results/BENCHMARK_LEDGER.md). Viết sau cùng, khi các chương kết quả đã xong.
\lipsumtodo{Tóm tắt tiểu luận.}


% ------------------------------------------------------------------ Nội dung (số Ả-rập, giữa cuối trang)
\clearpage
\pagenumbering{arabic}
\addchap{PHẦN MỞ ĐẦU}
% Theo mẫu: tính cấp thiết, mục đích, đối tượng và phạm vi, các công trình liên quan, kết quả dự kiến, bố cục.
% Khung đặt vấn đề (thống nhất 14/11/2026): bài toán tổng quát là "giải bài với bộ kiểm chứng một phần";
% Text-to-SQL là bài toán ứng dụng dùng làm ví dụ xuyên suốt; sinh mã nguồn là trường hợp được thực nghiệm.
% Nguồn ý: docs/report/KHUNG_BAO_CAO_v2.md mục 1-2; Rethink_extracted.md (mô hình, baseline, số rollout của bài báo).
% Phần này cố ý không nêu con số kết quả; số liệu nằm ở Chương 4-5 và Phần kết luận.

\addsec{TÍNH CẤP THIẾT CỦA ĐỀ TÀI}

Mô hình ngôn ngữ lớn (LLM) ngày càng được dùng để trực tiếp thực hiện công việc thay vì chỉ trả lời câu hỏi. Một ví dụ điển hình là trợ lý hỏi đáp trên dữ liệu doanh nghiệp: người dùng hỏi bằng ngôn ngữ tự nhiên, chẳng hạn ``doanh thu quý 3 của từng chi nhánh là bao nhiêu?'', hệ thống chuyển câu hỏi thành một truy vấn SQL, thực thi truy vấn trên cơ sở dữ liệu và trả lời từ kết quả. Bài toán này được gọi là Text-to-SQL \cite{yu2018spider,li2023bird}. Cái khó của nó nằm ở bước suy luận hơn là ở cú pháp. Một truy vấn có thể chạy không lỗi, trả về những con số trông hợp lý, nhưng vẫn sai vì nối nhầm bảng, quên loại các đơn hàng đã hủy hoặc lọc theo nhầm cột ngày. Khi LLM chỉ sinh một lần, những lỗi như vậy rất khó tự phát hiện.

Có hai cách chính để giảm loại lỗi này. Cách thứ nhất là dùng mô hình lớn hơn hoặc tinh chỉnh thêm, vốn đòi hỏi nhiều dữ liệu và tài nguyên huấn luyện. Cách thứ hai là giữ nguyên mô hình và dùng thêm tính toán \emph{lúc suy luận}: sinh nhiều lời giải, kiểm tra chúng, sửa dựa trên phản hồi, hoặc tìm kiếm có hệ thống trong không gian các hướng giải. Các nghiên cứu gần đây cho thấy tăng tính toán lúc suy luận có thể hiệu quả không kém tăng kích thước mô hình \cite{snell2024scaling}, và riêng việc lấy mẫu lặp lại đã nâng đáng kể tỷ lệ tìm được lời giải đúng khi có cách kiểm chứng tự động \cite{brown2024monkeys}.

Điều kiện ``có cách kiểm chứng tự động'' là then chốt. Trong Text-to-SQL, lúc trả lời, hệ thống chỉ kiểm tra được một phần: truy vấn có chạy được không, có đúng lược đồ dữ liệu không, có cho kết quả đúng trên một vài câu hỏi mẫu đã biết đáp án không. Còn truy vấn có đúng với ý người hỏi trên toàn bộ dữ liệu thật hay không thì lúc đó không biết được. Một truy vấn cho kết quả đúng trên một cơ sở dữ liệu vẫn có thể sai trên cơ sở dữ liệu khác \cite{zhong2020testsuite}. Có thể khái quát cấu trúc này thành một lớp bài toán chung, tạm gọi là \emph{giải bài với bộ kiểm chứng một phần}:

\begin{itemize}
  \item đầu vào là một yêu cầu $x$ bằng ngôn ngữ tự nhiên;
  \item LLM sinh ra một lời giải thực thi được $y$, ví dụ một truy vấn hoặc một chương trình;
  \item lúc giải chỉ có bộ kiểm chứng một phần $V_{\text{pub}}$: chấm được tự động, cho phản hồi khi sai, nhưng không đầy đủ;
  \item độ đúng thật sự do bộ kiểm chứng ẩn $V_{\text{hid}}$ quyết định, và bộ này không được dùng lúc giải;
  \item tổng chi phí gọi mô hình bị giới hạn bởi một ngân sách $B$.
\end{itemize}

Mục tiêu là trong ngân sách $B$, dùng $V_{\text{pub}}$ và phản hồi của nó để tìm lời giải $y$ có khả năng vượt qua $V_{\text{hid}}$ cao nhất. Phát biểu hình thức của bài toán được trình bày ở Chương~\ref{chap:lythuyet}.

Bài toán này đặc biệt đáng quan tâm với các \emph{LLM mở cỡ nhỏ}, từ 7 đến 9 tỷ tham số. Các mô hình này chạy được trên một GPU và triển khai được ngay trong hạ tầng của tổ chức, điều quan trọng khi dữ liệu doanh nghiệp không được phép gửi ra ngoài. Đổi lại, khả năng sinh đúng ngay lần đầu của chúng thấp hơn các mô hình lớn, nên tính toán thêm lúc suy luận là một cách bù đắp hấp dẫn. Câu hỏi thực tế đặt ra là: với cùng một ngân sách, một phương pháp tìm kiếm phức tạp có cho kết quả tốt hơn cách đơn giản là sinh nhiều lần rồi lọc bằng $V_{\text{pub}}$ hay không, và phải trả thêm bao nhiêu chi phí.

Để trả lời câu hỏi đó một cách khách quan, cần một môi trường mà $V_{\text{pub}}$ và $V_{\text{hid}}$ được tách sẵn và chấm được tự động. \emph{Sinh mã nguồn} từ đề bài thỏa đúng yêu cầu này: test công khai đi kèm đề bài đóng vai trò $V_{\text{pub}}$, bộ test ẩn đóng vai trò $V_{\text{hid}}$, và đã có các bộ dữ liệu chuẩn được dùng rộng rãi như APPS \cite{hendrycks2021apps} và HumanEval \cite{chen2021humaneval}. Sinh truy vấn SQL cũng là một dạng sinh mã. Vì vậy, tiểu luận chọn sinh mã nguồn làm bài toán thực nghiệm, còn Text-to-SQL là bài toán ứng dụng mà các kết quả hướng tới.

\addsec{MỤC ĐÍCH CỦA ĐỀ TÀI}

Tiểu luận nghiên cứu việc tăng cường khả năng suy luận của LLM mở cỡ nhỏ trong sinh mã nguồn bằng tìm kiếm cây Monte Carlo (MCTS) kết hợp phản hồi từ thực thi kiểm thử. Phương pháp được chọn để nghiên cứu là RethinkMCTS \cite{li2024rethinkmcts}. Phương pháp này tìm kiếm ở mức ý tưởng (thought) thay vì mức token, dùng phản hồi chi tiết theo từng khối mã để sửa lại những ý tưởng sai (rethink), và dùng chính LLM để đánh giá thêm các chương trình đã vượt qua test công khai (critic).

Cụ thể, tiểu luận cài đặt RethinkMCTS cùng hai phương pháp so sánh là sinh một lần (Direct) và sinh 16 lần rồi chọn theo test công khai (Best-of-16), chạy trên ba LLM mở cỡ nhỏ, và trả lời các câu hỏi nghiên cứu sau:

\begin{enumerate}
  \item RethinkMCTS có cho tỷ lệ giải đúng cao hơn Direct và Best-of-16 trên cả ba mô hình không?
  \item Mức cải thiện thay đổi thế nào giữa các mô hình, giữa mô hình chuyên mã nguồn và mô hình đa dụng, và giữa các mức độ khó của bài toán?
  \item Critic có giúp chọn được chương trình tốt hơn so với chỉ dựa vào test công khai không, và có thể dừng tìm kiếm sớm mà không giảm độ chính xác không?
  \item Chi phí tính toán (số lời gọi mô hình, số token) của từng phương pháp là bao nhiêu, và độ chính xác thay đổi thế nào theo ngân sách?
  \item Những sai lệch nhỏ nào trong cài đặt có thể làm mất tác dụng của các cơ chế tìm kiếm mà hệ thống không báo lỗi?
\end{enumerate}

Các câu hỏi trên được trả lời bằng thực nghiệm và kiểm định thống kê ở Chương~\ref{chap:ketqua} và Chương~\ref{chap:phantich}. Tiểu luận không giả định trước câu trả lời theo hướng nào.

\addsec{ĐỐI TƯỢNG VÀ PHẠM VI NGHIÊN CỨU}

\textbf{Đối tượng nghiên cứu} là các phương pháp tăng cường suy luận lúc suy luận cho LLM có sử dụng phản hồi từ thực thi, trọng tâm là RethinkMCTS, áp dụng cho bài toán sinh chương trình Python từ đề bài có test công khai và test ẩn.

\textbf{Phạm vi nghiên cứu:}
\begin{itemize}
  \item \emph{Mô hình:} ba LLM mở cỡ nhỏ thuộc ba thế hệ, gồm Qwen2.5-Coder-7B-Instruct (chuyên mã nguồn, 2024), Seed-Coder-8B-Instruct (chuyên mã nguồn, 2025) và Qwen3.5-9B (đa dụng, 2026, tắt chế độ suy nghĩ). Các mô hình được dùng nguyên trọng số, không tinh chỉnh.
  \item \emph{Dữ liệu:} APPS-300, gồm 300 bài lấy ngẫu nhiên có phân tầng từ tập test của APPS (100 bài mỗi mức khó), và toàn bộ 164 bài của HumanEval.
  \item \emph{Phương pháp so sánh:} Direct, Best-of-16 và RethinkMCTS với 16 rollout, tổng cộng 18 lần chạy (ba mô hình, hai bộ dữ liệu, ba phương pháp). Mỗi cấu hình chạy một lần.
  \item \emph{Hạ tầng:} suy luận trên GPU A100 40\,GB bằng vLLM \cite{kwon2023vllm}.
  \item \emph{Ngoài phạm vi:} tiểu luận không chạy thực nghiệm trên Text-to-SQL; bài toán này chỉ dùng làm bối cảnh ứng dụng và được bàn lại ở Phần kết luận. Các thí nghiệm loại bỏ từng thành phần (tắt rethink, thay phản hồi mức khối bằng thông tin thực thi đơn giản) chưa được thực hiện vì giới hạn kinh phí GPU; riêng vai trò của critic và của việc dừng sớm được phân tích lại từ dữ liệu đã chạy.
\end{itemize}

\addsec{CÁC CÔNG TRÌNH LIÊN QUAN}

\textbf{Lấy mẫu nhiều lần và lọc bằng test.} AlphaCode \cite{li2022alphacode} sinh rất nhiều chương trình cho mỗi bài rồi lọc bằng các ví dụ trong đề, cho thấy chỉ riêng lấy mẫu và lọc đã đạt mức thi đấu. CodeT \cite{chen2022codet} để LLM tự sinh thêm test rồi chọn chương trình theo mức đồng thuận khi thực thi. Brown và cộng sự \cite{brown2024monkeys} cho thấy tỷ lệ bài có ít nhất một lời giải đúng tăng đều theo số mẫu. Nhóm phương pháp này đơn giản, dễ song song hóa, và là mốc so sánh tự nhiên cho mọi phương pháp tìm kiếm phức tạp hơn.

\textbf{Sửa lỗi lặp dựa trên phản hồi thực thi.} Self-Debugging \cite{chen2023selfdebug} cho LLM đọc kết quả thực thi để tự sửa chương trình, và được đánh giá trên cả sinh mã lẫn Text-to-SQL. Reflexion \cite{shinn2023reflexion} lưu các nhận xét về lỗi trước đó vào bộ nhớ để dùng cho lần thử sau. LDB \cite{zhong2024ldb} chia chương trình thành các khối và theo dõi giá trị biến qua từng khối để định vị lỗi. Riêng với Text-to-SQL, việc thực thi các truy vấn trung gian để loại bỏ truy vấn lỗi đã được dùng từ trước khi có LLM \cite{wang2018egd}.

\textbf{Tìm kiếm có cấu trúc.} PG-TD \cite{zhang2023pgtd} dùng MCTS ở mức token với phần thưởng là tỷ lệ vượt qua test. Tree of Thoughts \cite{yao2023tot} tìm kiếm trên các bước suy nghĩ trung gian. LATS \cite{zhou2023lats} kết hợp MCTS với phản tư. RethinkMCTS \cite{li2024rethinkmcts} kết hợp các hướng trên: tìm kiếm ở mức ý tưởng, dùng phản hồi mức khối để sửa lại ý tưởng sai ngay trong cây, và kết hợp tỷ lệ vượt test với đánh giá của LLM làm phần thưởng. Phương pháp này được đánh giá trên APPS và HumanEval với GPT-3.5-turbo và GPT-4o-mini, so với các phương pháp trên ở cùng số rollout.

\textbf{Khoảng trống.} Các kết quả của RethinkMCTS mới được báo cáo trên hai mô hình đóng, mỗi cấu hình một lần chạy. Chưa rõ phương pháp còn giữ được lợi thế khi chuyển sang LLM mở cỡ nhỏ hay không, khi so với lấy mẫu nhiều lần ở cùng số chương trình sinh ra, có kiểm định thống kê ghép cặp trên cùng tập bài, và khi đặt độ chính xác cạnh chi phí tính toán. Đây là khoảng trống mà tiểu luận hướng tới.

\addsec{KẾT QUẢ DỰ KIẾN VÀ ĐÓNG GÓP}

\begin{enumerate}
  \item Một hệ thống cài đặt RethinkMCTS và các phương pháp so sánh, có kiểm thử, tách biệt nghiêm ngặt test công khai và test ẩn, lưu vết đầy đủ từng lời gọi mô hình và có thể chạy lại cho cùng kết quả.
  \item Một bộ kết quả thực nghiệm ghép cặp trên ba LLM mở cỡ nhỏ và hai bộ dữ liệu, với kiểm định McNemar và hiệu chỉnh đa so sánh Holm.
  \item Phân tích quan hệ giữa độ chính xác và chi phí, vai trò của critic và của việc dừng sớm, cùng các bài học về những sai lệch cài đặt có thể làm mất tác dụng của cơ chế tìm kiếm.
  \item Các nhận xét về điều kiện để áp dụng hướng tiếp cận này cho bài toán ứng dụng Text-to-SQL.
\end{enumerate}

\addsec{BỐ CỤC TIỂU LUẬN}

Ngoài Phần mở đầu và Phần kết luận, tiểu luận gồm năm chương.
\begin{itemize}
  \item \textbf{Chương~\ref{chap:lythuyet}} trình bày cơ sở lý thuyết: sinh mã nguồn bằng LLM, phát biểu bài toán giải bài với bộ kiểm chứng một phần, các phương pháp tìm kiếm lúc suy luận, RethinkMCTS và các kiểm định thống kê được sử dụng.
  \item \textbf{Chương~\ref{chap:caidat}} mô tả thiết kế và cài đặt hệ thống.
  \item \textbf{Chương~\ref{chap:thietlap}} mô tả thiết lập thực nghiệm: dữ liệu, mô hình, cấu hình, phương pháp so sánh và độ đo.
  \item \textbf{Chương~\ref{chap:ketqua}} trình bày kết quả thực nghiệm.
  \item \textbf{Chương~\ref{chap:phantich}} phân tích vai trò của từng thành phần, hành vi tìm kiếm, bài học từ quá trình cài đặt và các yếu tố ảnh hưởng đến tính hợp lệ của kết quả.
\end{itemize}
Phần kết luận tóm tắt kết quả, đóng góp, hạn chế và hướng phát triển, trong đó có việc áp dụng cho Text-to-SQL.


\phantoc{PHẦN NỘI DUNG}
\chapter{CƠ SỞ LÝ THUYẾT}
\label{chap:lythuyet}

\section{SINH MÃ NGUỒN BẰNG MÔ HÌNH NGÔN NGỮ LỚN}
\lipsumtodo{LLM cho code, đánh giá bằng test, Pass@1; các bộ dữ liệu APPS \cite{hendrycks2021apps} và HumanEval \cite{chen2021humaneval}.
Quy ước thuật ngữ: dùng ``LLM'' làm thuật ngữ chung (khớp tên đề tài); ba mô hình của tiểu luận gọi là ``LLM mở cỡ nhỏ (7--9B tham số)''. Nêu một câu rằng một số tài liệu gọi nhóm mô hình cỡ này là small language model (SLM), nhưng ngưỡng SLM không thống nhất, nên tiểu luận không dựa vào ngưỡng nào và luôn ghi số tham số cụ thể.}

\section{BÀI TOÁN GIẢI BÀI VỚI BỘ KIỂM CHỨNG MỘT PHẦN}
\label{sec:baitoan}
\lipsumtodo{Phát biểu hình thức: yêu cầu $x$, lời giải $y$, bộ kiểm chứng một phần $V_{\text{pub}}$ (có phản hồi), bộ kiểm chứng ẩn $V_{\text{hid}}$, ngân sách $B$; mục tiêu tối đa xác suất $y$ vượt $V_{\text{hid}}$. Bảng ánh xạ hai trường hợp: sinh mã (public test / hidden test) và Text-to-SQL (chạy được, đúng lược đồ, đúng trên câu hỏi mẫu / đúng trên dữ liệu thật). Nêu hệ quả: khi nhiều lời giải cùng vượt $V_{\text{pub}}$, cần thêm tín hiệu để phân xử (critic).}

\section{TÌM KIẾM LÚC SUY LUẬN (TEST-TIME SEARCH)}
\subsection{Lấy mẫu nhiều lần và Best-of-N}
\lipsumtodo{}
\subsection{Monte Carlo Tree Search và công thức P-UCB}
\lipsumtodo{Selection, expansion, evaluation, backup (Max); P-UCB (Eq.~1--2 của bài báo).}

\section{RETHINKMCTS}
\subsection{Tìm kiếm ở mức thought}
\lipsumtodo{}
\subsection{Rethink tại lá và phản hồi mức block}
\lipsumtodo{CFG, trace biến, verbal feedback.}
\subsection{Dual reward có điều kiện}
\lipsumtodo{Critic chỉ chạy khi mọi public test đều pass (Eq.~3), thang $[-1,1]$.}

\section{KIỂM ĐỊNH THỐNG KÊ CHO SO SÁNH GHÉP CẶP}
\lipsumtodo{McNemar exact \cite{mcnemar1947}, hiệu chỉnh Holm \cite{holm1979}, khoảng tin cậy Wilson \cite{wilson1927}, bootstrap ghép cặp.}

\chapter{THIẾT KẾ VÀ CÀI ĐẶT HỆ THỐNG}
\label{chap:caidat}

\section{YÊU CẦU VÀ NGUYÊN TẮC THIẾT KẾ}
\lipsumtodo{Tái hiện đúng bài báo, chống rò rỉ hidden test, mọi run có thể kiểm tra lại (manifest, hash).}

\section{KIẾN TRÚC TỔNG THỂ}
\lipsumtodo{Ports \& adapters: domain, policies, ports, adapters, application; bootstrap là nơi duy nhất nối các thành phần. Sơ đồ thành phần/lớp (UML).}

\section{CÁC THÀNH PHẦN CHÍNH}
\subsection{Bộ máy tìm kiếm MCTS}
\lipsumtodo{}
\subsection{Đánh giá ứng viên và phản hồi}
\lipsumtodo{CandidateEvaluator, FeedbackService, phân tích block bằng CFG và trace.}
\subsection{Thực thi mã và chấm điểm}
\lipsumtodo{Executor, HiddenJudge, sandbox.}
\subsection{Lưu vết và tiếp tục chạy}
\lipsumtodo{Artifact schema 2, manifest, resume theo từng bài.}

\section{TRIỂN KHAI TRÊN GPU}
\lipsumtodo{vLLM, Colab A100 và Modal; quy trình đóng băng commit.}

\chapter{THIẾT LẬP THỰC NGHIỆM}
\label{chap:thietlap}
% Nguồn (chỉ để truy vết, không in ra): docs/research/PROTOCOL.md mục 8, configs/v3_7b/README.md,
% configs/v3_7b/apps_mcts.yaml, docs/research/DECISION_LOG.md (D-006 đến D-010),
% docs/results/BENCHMARK_LEDGER.md (mục 1, 11, 12), data/cohorts/apps_benchmark_300.json.

Chương này mô tả toàn bộ thiết lập dùng cho các thực nghiệm ở Chương~\ref{chap:ketqua} và Chương~\ref{chap:phantich}: dữ liệu, mô hình, phần cứng, cấu hình thuật toán, các phương pháp so sánh, độ đo và quy trình chạy. Mọi lần chạy được báo cáo trong tiểu luận đều theo cùng một bộ quy tắc, được cố định trước khi chạy dữ liệu chính thức, để không có tham số nào được chỉnh sau khi đã nhìn thấy kết quả.

\section{DỮ LIỆU VÀ CÁCH CHIA TEST}
\label{sec:dulieu}

Tiểu luận dùng hai bộ dữ liệu sinh mã nguồn Python: APPS \cite{hendrycks2021apps} và HumanEval \cite{chen2021humaneval}.

\subsection{APPS-300}

APPS gồm các bài lập trình thi đấu; chương trình đọc dữ liệu từ đầu vào chuẩn và ghi kết quả ra đầu ra chuẩn. Tập test của APPS có 5.000 bài, chia thành ba mức khó: introductory, interview và competition. Do chi phí chạy MCTS trên toàn bộ tập là quá lớn, tiểu luận lấy một tập con cố định 300 bài, gọi là \emph{APPS-300}:

\begin{itemize}
  \item Chỉ xét các bài có ít nhất 3 test và có số đầu vào bằng số đầu ra. Số bài thỏa điều kiện là 758 (introductory), 2.712 (interview) và 640 (competition).
  \item Lấy mẫu ngẫu nhiên phân tầng, không hoàn lại, mỗi mức khó 100 bài, với seed 2027.
  \item Danh sách 300 bài được cố định và dùng chung cho mọi mô hình và mọi phương pháp. Không có bước lọc nào khác sau đó.
\end{itemize}

Bộ test của mỗi bài được chia thành hai phần: \emph{public test} để thuật toán dùng trong quá trình tìm kiếm, và \emph{hidden test} để chấm kết quả cuối cùng. Số public test bằng phần nguyên của một nửa số test, phần còn lại là hidden. Phép chia được thực hiện trên các test \emph{khác nhau}, vì một số bài trong APPS liệt kê lặp lại cùng một test; chẳng hạn \texttt{apps\_4349} có 10 test nhưng chỉ 5 test khác nhau. Nếu chia trên danh sách gốc, một test có thể vừa nằm trong public vừa nằm trong hidden, khiến hidden không còn là dữ liệu chưa thấy. Mỗi test có giới hạn thời gian 4 giây.

Bảng~\ref{tab:phanbotest} cho thấy số test rất khác nhau giữa các mức khó. Ở mức competition, trung vị chỉ là 2 public test và 2 hidden test; 48 bài chỉ có đúng một public test và 57 bài có không quá hai hidden test. Với những bài như vậy, một chương trình sai vẫn có thể vượt qua toàn bộ hidden test. Điểm này được thảo luận lại ở Mục~\ref{sec:hople}.

\begin{table}[htbp]
  \caption{Phân bố số public test và hidden test trong APPS-300 theo mức khó}
  \label{tab:phanbotest}
  % Nguồn: prepared/apps_half (public/problems.jsonl, judge/suites) + data/cohorts/apps_benchmark_300.json.
  \begin{tabularx}{\linewidth}{>{\raggedright\arraybackslash}p{2.6cm}*{4}{>{\centering\arraybackslash}X}}
    \toprule
    Mức khó & Số public test: trung vị (min--max) & Số hidden test: trung vị (min--max) & Số bài chỉ có 1 public test & Số bài có $\le 2$ hidden test \\
    \midrule
    Introductory & 4 (1--67)    & 5 (2--67)  & 24 & 34 \\
    Interview    & 10,5 (1--62) & 11 (2--62) & 13 & 16 \\
    Competition  & 2 (1--56)    & 2 (2--57)  & 48 & 57 \\
    \bottomrule
  \end{tabularx}
  \ghichubang{Mỗi mức khó gồm 100 bài.}
\end{table}

\subsection{HumanEval-164}

HumanEval gồm 164 bài viết hàm Python theo mô tả trong docstring. Public test là các ví dụ đi kèm đề bài (trung vị 2 ví dụ mỗi bài, từ 1 đến 8). Hidden test là hàm kiểm tra \texttt{check} chính thức của bộ dữ liệu. Toàn bộ hàm kiểm tra được chấm như một khối: chương trình chỉ được tính là đúng khi mọi \texttt{assert} trong đó đều đúng. Mỗi lần chạy kiểm tra có giới hạn 3 giây. Tiểu luận dùng toàn bộ 164 bài.

\subsection{Tách biệt public và hidden test}

Để tránh rò rỉ hidden test vào quá trình tìm kiếm, hệ thống tách dữ liệu thành hai đối tượng riêng (Chương~\ref{chap:caidat}). Mọi phương pháp sinh mã chỉ nhận đề bài và public test. Hidden test chỉ được nạp bởi bộ chấm, và chỉ sau khi chương trình được chọn của bài đó đã được ghi ra đĩa. Như vậy, việc chọn chương trình cuối cùng không thể phụ thuộc vào kết quả hidden.

\section{MÔ HÌNH VÀ PHẦN CỨNG}
\label{sec:mohinh}

\subsection{Mô hình}

Tiểu luận dùng ba LLM mở cỡ nhỏ, từ 7 đến 9 tỷ tham số (Bảng~\ref{tab:mohinh}). Hai mô hình đầu là mô hình chuyên cho mã nguồn; mô hình thứ ba là mô hình đa dụng. Cả ba được dùng nguyên trọng số, không tinh chỉnh.

\begin{table}[htbp]
  \caption{Các mô hình được sử dụng}
  \label{tab:mohinh}
  \begin{tabular}{l>{\footnotesize}ll}
    \toprule
    Tên trong tiểu luận & \normalsize Định danh trên Hugging Face & Loại \\
    \midrule
    Qwen2.5-Coder-7B & \texttt{Qwen/Qwen2.5-Coder-7B-Instruct} & Chuyên mã nguồn \\
    Seed-Coder-8B & \texttt{ByteDance-Seed/Seed-Coder-8B-Instruct} & Chuyên mã nguồn \\
    Qwen3.5-9B & \texttt{Qwen/Qwen3.5-9B} & Đa dụng \\
    \bottomrule
  \end{tabular}
\end{table}

Qwen2.5-Coder-7B và Seed-Coder-8B không có chế độ suy nghĩ (thinking). Qwen3.5-9B mặc định suy nghĩ dài trước khi trả lời; chế độ này được tắt cho mọi phương pháp bằng tham số \texttt{enable\_thinking: false} của chat template, để ba mô hình được so sánh trong cùng điều kiện: mô hình trả lời trực tiếp, mọi quá trình suy luận thêm đều đến từ thuật toán tìm kiếm.

\subsection{Phần cứng và phần mềm}

Mô hình được phục vụ bằng vLLM \cite{kwon2023vllm} phiên bản 0.31.0, độ chính xác bf16, độ dài ngữ cảnh tối đa 32.768 token, seed 0 và có bật prefix caching. vLLM và chương trình tìm kiếm chạy chung trên một máy có GPU NVIDIA A100-SXM4-40GB. Chương trình tìm kiếm, phần trace và phần chạy test dùng Python 3.11. Phiên bản này được cố định vì ở Python 3.13, ba kiểm thử về trace bị lỗi (hàm đệ quy bị bỏ sót lượt chạy), làm phản hồi mức block thiếu dữ liệu mà không báo lỗi. Mỗi lần chạy giải 16 bài song song; mỗi bài chạy tối đa 4 test cùng lúc và toàn hệ thống có tối đa 8 tiến trình test.

Các lần chạy được thực hiện trên hai nền tảng thuê GPU theo giờ là Google Colab và Modal, cùng loại GPU và cùng gói mã nguồn:

\begin{itemize}
  \item Qwen2.5-Coder-7B và Seed-Coder-8B: cả 12 lần chạy trên Colab.
  \item Qwen3.5-9B: lần chạy RethinkMCTS trên APPS-300 dùng Colab; năm lần chạy còn lại dùng Modal. Lý do là trong ngày 13/11/2026, máy Colab liên tục bị thu hồi sau khoảng 23--29 phút, không đủ cho một lần chạy MCTS dài. Trước khi chạy, bước kiểm tra cấu hình (gate, Mục~\ref{sec:quytrinh}) được chạy lại trên Modal và cho kết quả khớp với Colab ở 12/13 bài.
\end{itemize}

Việc ô APPS của Qwen3.5-9B chạy trên hai nền tảng là một sai lệch so với quy tắc ban đầu (mọi phương pháp của một mô hình chạy trên cùng nền tảng), và được ghi nhận ở Mục~\ref{sec:hople}.

\section{CẤU HÌNH THUẬT TOÁN}
\label{sec:cauhinh}

Bảng~\ref{tab:cauhinh} liệt kê các tham số của RethinkMCTS dùng trong mọi lần chạy. Ý nghĩa của từng tham số được trình bày ở Chương~\ref{chap:lythuyet}.

\begin{table}[htbp]
  \caption{Cấu hình RethinkMCTS}
  \label{tab:cauhinh}
  % Nguồn: configs/v3_7b/apps_mcts.yaml, configs/v3_7b/README.md, docs/research/PROTOCOL.md mục 8.
  \begin{tabularx}{\linewidth}{>{\raggedright\arraybackslash}p{5cm}>{\raggedright\arraybackslash}X}
    \toprule
    Tham số & Giá trị \\
    \midrule
    Ngân sách tìm kiếm & 16 rollout; lượt rethink không tính vào ngân sách \\
    Độ rộng mở rộng & 3 thought con mỗi node \\
    Số lần rethink mỗi node & 1 \\
    Chọn node & P-UCB, $c = 4$, $c_{\text{base}} = 10$ \\
    Lan truyền giá trị & Max \\
    Phần thưởng & $0{,}8\,v_{\text{test}} + 0{,}2\,v_{\text{llm}}$; critic cho điểm trong $[-1, 1]$, chỉ được gọi khi chương trình pass toàn bộ public test \\
    Phản hồi khi chương trình sai & Phản hồi mức block (phân tích theo từng khối lệnh) \\
    Dừng sớm & Không; phương án dừng sớm được tính lại offline \\
    Chọn kết quả & Chương trình có phần thưởng lớn nhất \\
    Temperature & 0 \\
    Giới hạn an toàn mỗi bài & 4 triệu token, 400 lời gọi, 3.600 giây \\
    \bottomrule
  \end{tabularx}
\end{table}

Mỗi vai trò gọi mô hình có giới hạn số token đầu ra riêng: sinh thought (expansion) 1.024, sinh mã 4.096, rethink 512, critic 256 và phân tích block 3.072. Các giới hạn này được đặt sau khi lần chạy thử đầu tiên với giới hạn nhỏ hơn cho thấy 5,9\% lời gọi của Qwen3.5-9B bị cắt giữa chừng, và được áp dụng chung cho cả ba mô hình.

\section{CÁC PHƯƠNG PHÁP SO SÁNH}
\label{sec:phuongphap}

Ba phương pháp được chạy trên mỗi mô hình và mỗi bộ dữ liệu, tổng cộng $3 \times 3 \times 2 = 18$ lần chạy. Cả ba dùng chung prompt sinh mã.

\begin{description}[style=nextline,leftmargin=1cm]
  \item[Direct] Sinh một chương trình duy nhất với temperature 0. Đây là mức năng lực của mô hình khi không có tìm kiếm.
  \item[Best-of-16] Sinh 16 chương trình độc lập với temperature 0,8 và top-p 0,95, chạy cả 16 trên public test, chọn chương trình có tỷ lệ pass public cao nhất; nếu bằng nhau thì lấy chương trình sinh trước. Đây là baseline lấy mẫu với cùng số lần sinh mã ban đầu như RethinkMCTS (16 rollout).
  \item[RethinkMCTS] Thuật toán tìm kiếm với cấu hình ở Bảng~\ref{tab:cauhinh}.
\end{description}

Best-of-16 và RethinkMCTS có cùng số rollout nhưng không cùng chi phí: RethinkMCTS còn gọi mô hình cho các bước sinh thought, rethink, phân tích block và critic. Vì vậy, tiểu luận báo cáo chi phí (token và số lời gọi) song song với độ chính xác, và so sánh thêm với Best-of-$k$ ($k = 1, \ldots, 16$) tính lại từ 16 mẫu đã sinh (Chương~\ref{chap:ketqua}).

Hai thí nghiệm ablation trên Qwen2.5-Coder-7B (tắt rethink; chỉ dùng thông tin thực thi thay vì phản hồi mức block) đã được lên kế hoạch nhưng chưa được thực hiện trong phạm vi tiểu luận vì giới hạn kinh phí GPU. Riêng vai trò của critic và của việc dừng sớm được phân tích offline từ các lần chạy RethinkMCTS sẵn có, không cần chạy thêm.

\section{ĐỘ ĐO VÀ KIỂM ĐỊNH}
\label{sec:dodo}

\subsection{Độ đo}

\begin{itemize}
  \item \textbf{Pass@1 (độ đo chính):} tỷ lệ bài mà chương trình được chọn vượt qua \emph{toàn bộ} hidden test. Mẫu số là toàn bộ bài của bộ dữ liệu; bài lỗi kỹ thuật hoặc không sinh được chương trình nào được tính là không pass.
  \item \textbf{Pass rate:} trung bình theo bài của tỷ lệ hidden test mà chương trình được chọn vượt qua.
  \item \textbf{Chi phí:} số token đầu vào và đầu ra, số lời gọi mô hình (tính cả expansion, sinh mã, rethink, critic và phân tích block) và thời gian giải mỗi bài. Lời gọi được thử lại do lỗi mạng được đếm riêng; lời gọi không có thông tin token được ghi là thiếu, không ghi là 0.
\end{itemize}

\subsection{Kiểm định thống kê}

Vì mọi phương pháp chạy trên cùng các bài, so sánh giữa hai phương pháp là so sánh \emph{ghép cặp} theo từng bài:

\begin{itemize}
  \item Khoảng tin cậy 95\% của Pass@1 tính theo phương pháp Wilson \cite{wilson1927}.
  \item Kiểm định McNemar dạng chính xác \cite{mcnemar1947} trên số bài mà chỉ một trong hai phương pháp giải được.
  \item Có 18 phép so sánh (3 mô hình $\times$ 2 bộ dữ liệu $\times$ 3 cặp phương pháp), nên p-value được hiệu chỉnh bằng phương pháp Holm \cite{holm1979} trên cả 18 phép. Một khác biệt chỉ được xem là có ý nghĩa thống kê khi p-value sau hiệu chỉnh nhỏ hơn 0,05.
  \item Khoảng tin cậy 95\% của chênh lệch Pass@1 tính bằng bootstrap ghép cặp theo bài, 10.000 lần lấy mẫu lại.
\end{itemize}
% Tính bằng scripts/analyze_v3.py (bootstrap seed 0); kết quả: docs/results/v3/analysis_v3_3models.json.

\section{QUY TRÌNH THỰC NGHIỆM}
\label{sec:quytrinh}

Quy trình gồm bốn bước, theo thứ tự:

\begin{enumerate}[label=\arabic*.]
  \item \textbf{Pilot.} Chạy Qwen2.5-Coder-7B trên 20 bài APPS \emph{nằm ngoài} APPS-300, với cả ba phương pháp. Pilot dùng để kiểm tra rằng mọi thành phần hoạt động đúng và ước lượng chi phí. Tiêu chí đạt gồm: không có lỗi template prompt; tối đa 1/20 bài lỗi kỹ thuật; ít nhất 90\% kết quả phân tích block được đưa vào phản hồi; đủ 16 rollout và có cây đạt độ sâu từ 3 trở lên; không quá 5\% lời gọi bị cắt; không bài nào chạm giới hạn 3.600 giây. Kết quả pilot không được dùng làm kết quả của tiểu luận.
  \item \textbf{Cố định cấu hình.} Sau khi pilot đạt, mã nguồn, prompt và cấu hình được cố định (ngày 12/11/2026). Mọi lần chạy chính thức dùng đúng phiên bản này. Không được đổi cấu hình, prompt, mã nguồn hay danh sách bài sau khi đã bắt đầu chạy; nếu buộc phải sửa thì bỏ toàn bộ kết quả của mô hình đó và chạy lại.
  \item \textbf{Kiểm tra trước cho từng mô hình (gate).} Trước khi chạy một mô hình, chạy thử trên 5 bài APPS ngoài APPS-300 (RethinkMCTS và Direct) và 3 bài HumanEval (chỉ để kiểm tra định dạng đầu ra). Nếu không đạt thì dừng, không chạy dữ liệu chính thức.
  \item \textbf{Chạy chính thức.} Chạy 18 lần, mỗi tổ hợp mô hình $\times$ bộ dữ liệu $\times$ phương pháp ghi vào một thư mục riêng, kèm một file manifest ghi lại cấu hình, mã băm của dữ liệu, mã nguồn và prompt.
\end{enumerate}
% Commit đóng băng: d9666c4.

Trong đợt chạy, máy Colab nhiều lần bị thu hồi giữa chừng. Hệ thống xử lý bằng cách chạy tiếp theo từng bài: bài đã có kết quả được giữ nguyên, bài đang giải dở thì được giải lại từ đầu, và việc chạy tiếp chỉ được phép khi manifest khớp hoàn toàn với lần chạy trước. Có 3 bài Best-of-16 của Qwen2.5-Coder-7B bị một lỗi kiểm tra khi chạy tiếp (do ký tự xuống dòng kiểu Windows trong mã sinh ra) nên được giải lại; tiêu chí chọn các bài này không phụ thuộc vào kết quả chấm, nên không gây thiên lệch.
% Chi tiết vận hành: docs/results/BENCHMARK_LEDGER.md mục 12.

Kết thúc đợt chạy, cả 18 lần chạy đều hoàn tất đủ số bài, không có bài nào lỗi kỹ thuật. Chấm lại hidden test cho chương trình được RethinkMCTS chọn trên máy cá nhân cho cùng kết quả với máy chạy ở 298--300/300 bài APPS và 164/164 bài HumanEval, tùy mô hình.

\chapter{KẾT QUẢ THỰC NGHIỆM}
\label{chap:ketqua}
% Mọi con số lấy từ docs/results/BENCHMARK_LEDGER.md và docs/results/v3/analysis_v3_3models.json.
% Hình sinh bằng scripts/make_figures_v3.py (thư mục docs/report/figures/v3/).

\section{ĐỘ PHỦ VÀ LỖI KỸ THUẬT}
\lipsumtodo{18 run, không có lỗi kỹ thuật; các bài không có candidate được tính 0.}

\section{PASS@1 TRÊN APPS-300 VÀ HUMANEVAL-164}
\lipsumtodo{Bảng B4.2, B4.4.}

\begin{figure}[htbp]
  \includegraphics[width=\linewidth]{H1_pass_at_1.pdf}
  \caption{Pass@1 trên hidden test của ba phương pháp, kèm khoảng tin cậy Wilson 95\%}
  \label{fig:pass1}
\end{figure}

\section{SO SÁNH GHÉP CẶP}
\lipsumtodo{Bảng B4.3.}

\begin{figure}[htbp]
  \includegraphics[width=\linewidth]{H2_paired_forest.pdf}
  \caption{Chênh lệch Pass@1 ghép cặp và ý nghĩa thống kê sau hiệu chỉnh Holm}
  \label{fig:forest}
\end{figure}

\section{KẾT QUẢ THEO MỨC KHÓ TRÊN APPS}
\begin{figure}[htbp]
  \includegraphics[width=\linewidth]{H4_apps_by_difficulty.pdf}
  \caption{Pass@1 trên APPS-300 theo mức khó}
  \label{fig:dokho}
\end{figure}
\lipsumtodo{}

\section{CHI PHÍ VÀ HIỆU NĂNG THEO NGÂN SÁCH}
\begin{figure}[htbp]
  \includegraphics[width=\linewidth]{H3_cost_vs_pass.pdf}
  \caption{Pass@1 theo số token trung bình mỗi bài: Best-of-$k$, Direct và RethinkMCTS}
  \label{fig:chiphi}
\end{figure}
\begin{figure}[htbp]
  \includegraphics[width=\linewidth]{H6_mcts_tokens_by_role.pdf}
  \caption{Token trung bình mỗi bài của RethinkMCTS theo vai trò}
  \label{fig:vaitro}
\end{figure}
\lipsumtodo{Bảng B4.5.}

\chapter{PHÂN TÍCH VÀ THẢO LUẬN}
\label{chap:phantich}

\section{VAI TRÒ CỦA TỪNG THÀNH PHẦN TRONG MCTS}
\begin{figure}[htbp]
  \includegraphics[width=\linewidth]{H5_mcts_stages.pdf}
  \caption{Pass@1 của RethinkMCTS qua từng bước: chương trình đầu tiên, sau search, sau critic}
  \label{fig:buoc}
\end{figure}
\lipsumtodo{Search so với chương trình đầu tiên; critic chọn lại offline (B5.2).}

\section{HÀNH VI TÌM KIẾM}
\lipsumtodo{Độ sâu cây, số ứng viên, nguồn gốc lời giải đúng, dừng sớm tính lại offline (B5.4).}

\section{BÀI HỌC TỪ QUÁ TRÌNH TÁI HIỆN}
\lipsumtodo{Các lệch của cài đặt cũ và cách v3 sửa (B5.3, D-006); mô tả, không dùng số liệu v1.}

\section{PHÂN LOẠI LỖI VÀ NGHIÊN CỨU TÌNH HUỐNG}
\lipsumtodo{Chỉ chọn case từ run v3.}

\section{CÁC YẾU TỐ ẢNH HƯỞNG ĐẾN TÍNH HỢP LỆ}
\label{sec:hople}
\lipsumtodo{Một seed; nhiễm dữ liệu; temperature 0 trên vLLM không tất định tuyệt đối; Qwen3.5-9B bị cắt output trên APPS và chạy trên hai nền tảng (D-010); Competition có ít hidden test.}


\addchap{PHẦN KẾT LUẬN}
% Theo mẫu: ngắn gọn, không bàn luận thêm.

\addsec{KẾT QUẢ ĐẠT ĐƯỢC}
\lipsumtodo{}

\addsec{ĐÓNG GÓP MỚI}
\lipsumtodo{}

\addsec{ƯU ĐIỂM VÀ NHƯỢC ĐIỂM}
\lipsumtodo{}

\addsec{HƯỚNG PHÁT TRIỂN}
\lipsumtodo{Quay lại bài toán ứng dụng Text-to-SQL đã nêu ở Phần mở đầu: áp dụng cho các bước kiểm chứng được trong trợ lý hỏi đáp dữ liệu (truy vấn chạy được, đúng lược đồ, đúng trên câu hỏi mẫu). Điều kiện rút ra từ kết quả: hiệu quả phụ thuộc chất lượng của bộ kiểm chứng; LLM tự đánh giá (critic) chưa đủ tin cậy để thay thế. Các thí nghiệm loại bỏ thành phần chưa chạy.}


\printbibliography[heading=tailieu]

\appendix
\chapter{PROMPT, CẤU HÌNH VÀ MANIFEST}
\label{app:cauhinh}
\lipsumtodo{Prompt, config và manifest hash của các run v3.1.}

\chapter{KẾT QUẢ PILOT}
\label{app:pilot}
\lipsumtodo{Chỉ để minh bạch quy trình, không phải kết quả.}


\end{document}
